\begin{center}
	{\Large \textbf{Appendix A}}-\vspace*{0.5 cm}
	{\Large \textbf{Key Source Code Listings}}
	
\end{center}

\paragraph{Listing A.1: ECGNet Model Architecture (\texttt{inference/heart.py})}
\begin{lstlisting}[language=Python, caption={ECGNet Multi-Scale Attention CNN Architecture}]
class ChannelAttention(nn.Module):
    def __init__(self, channels, r=8):
        super().__init__()
        self.avg = nn.AdaptiveAvgPool1d(1)
        self.fc = nn.Sequential(
            nn.Linear(channels, channels // r),
            nn.ReLU(),
            nn.Linear(channels // r, channels),
            nn.Sigmoid()
        )

    def forward(self, x):
        b, c, t = x.size()
        y = self.avg(x).view(b, c)
        y = self.fc(y).view(b, c, 1)
        return x * y


class MultiScaleBlock(nn.Module):
    def __init__(self, in_ch, out_ch):
        super().__init__()
        self.conv3 = nn.Conv1d(in_ch, out_ch, 3, padding=1)
        self.conv5 = nn.Conv1d(in_ch, out_ch, 5, padding=2)
        self.conv7 = nn.Conv1d(in_ch, out_ch, 7, padding=3)
        self.bn = nn.BatchNorm1d(out_ch * 3)
        self.att = ChannelAttention(out_ch * 3)

    def forward(self, x):
        x = torch.cat([self.conv3(x), self.conv5(x),
                       self.conv7(x)], dim=1)
        x = F.relu(self.bn(x))
        return self.att(x)


class ECGNet(nn.Module):
    def __init__(self):
        super().__init__()
        self.block1 = MultiScaleBlock(1, 32)
        self.pool1  = nn.MaxPool1d(2)
        self.block2 = MultiScaleBlock(96, 64)
        self.pool2  = nn.MaxPool1d(2)
        self.block3 = MultiScaleBlock(192, 128)
        self.gap    = nn.AdaptiveAvgPool1d(1)
        self.fc     = nn.Linear(384, 5)

    def forward(self, x):
        x = self.pool1(self.block1(x))
        x = self.pool2(self.block2(x))
        x = self.block3(x)
        x = self.gap(x).squeeze(-1)
        return self.fc(x)
\end{lstlisting}

\paragraph{Listing A.2: Triage Logic (\texttt{chat/parser.py})}
\begin{lstlisting}[language=Python, caption={Automated Triage Classification Function}]
def determine_triage(cardiac_risk, metabolic_risk, motor_risk):
    high_levels     = {'high', 'elevated'}
    moderate_levels = {'moderate', 'borderline', 'mild'}
    risks = [cardiac_risk.lower(), metabolic_risk.lower(),
             motor_risk.lower()]

    if any(r in high_levels for r in risks):
        return "priority_review"
    moderate_count = sum(1 for r in risks if r in moderate_levels)
    if moderate_count >= 2:
        return "recommended_check"
    return "routine"
\end{lstlisting}

\paragraph{Listing A.3: DiabetesNet Architecture (\texttt{inference/diabetes.py})}
\begin{lstlisting}[language=Python, caption={DiabetesNet Tabular Deep Neural Network}]
class DiabetesNet(nn.Module):
    def __init__(self):
        super().__init__()
        self.net = nn.Sequential(
            nn.Linear(8, 32),
            nn.BatchNorm1d(32),
            nn.ReLU(),
            nn.Dropout(0.3),
            nn.Linear(32, 16),
            nn.ReLU(),
            nn.Dropout(0.2),
            nn.Linear(16, 1)
        )

    def forward(self, x):
        return self.net(x)
\end{lstlisting}

